\section{Step 111 Admit the clip before opening files}\label{h263-clip-s01}
Check aligned request/result records and all non-wrapping extents. Admit bounded
picture count, clock step, quantizer and mode; then evaluate the duration product.
Scan each bounded path, reject empty or unterminated names and compare all three
pairs for equality. Leave progress untouched on rejection. Only after these checks
initialize the accepted count, byte total and reference pointer to zero.

\paragraph{Observe} Two different path pointers can still contain the same filename.
\Needspace{8\baselineskip}
\paragraph{Check} What prevents an output open from truncating the input?
\begin{itemize}
  \item[A.] Comparing admitted path contents before opening files.
  \item[B.] Assuming different pointer values imply different names.
  \item[C.] Choosing descriptor three in advance.
\end{itemize}

\section{Step 112 Own each successful descriptor}\label{h263-clip-s02}
Initialize owned-descriptor fields to a negative sentinel. Open the input for
reading, then both outputs for writing, storing each actual returned descriptor.
On an open failure, enter cleanup with only the successfully acquired descriptors.
The same cleanup path handles later transfer or child failures. Preserve the
intended return status across close syscalls.

\paragraph{Observe} If the third open fails, the first two descriptors still require closing.
\Needspace{8\baselineskip}
\paragraph{Check} Which descriptor is owned?
\begin{itemize}
  \item[A.] Every integer from zero through thirty-one.
  \item[B.] Only a nonnegative descriptor actually returned by this call's successful open.
  \item[C.] The input path pointer.
\end{itemize}

\section{Step 113 Load exactly one raw picture}\label{h263-clip-s03}
Read at most 4096 bytes per request into the reusable raw buffer. Advance the
cursor and reduce the remaining count by the returned positive count. Continue
until all 38\,016 bytes are present. On early EOF or a negative read, report input
failure without invoking the picture encoder for this partial buffer.

\paragraph{Observe} A read returning 137 bytes accepts exactly 137 even when 4096 were requested.
\Needspace{8\baselineskip}
\paragraph{Check} What advances the destination cursor?
\begin{itemize}
  \item[A.] Always 4096 bytes.
  \item[B.] The complete picture size on the first call.
  \item[C.] The actual positive return count.
\end{itemize}

\section{Step 114 Construct the next picture request}\label{h263-clip-s04}
Use no reference for picture zero and the accepted reconstructed reference for
later pictures. Select the other reconstruction buffer as destination. Reset the
private picture context to zero used bits, supply the fixed workspace and encoded
buffer, and pass quantizer/mode. Compute temporal reference as picture index times
clock step modulo 256. Use GFID zero for the initial I and one for later P pictures.

\paragraph{Observe} Picture index two at step six has temporal reference twelve, regardless of output byte length.
\Needspace{8\baselineskip}
\paragraph{Check} Which reference enters a P-picture?
\begin{itemize}
  \item[A.] The last accepted reconstruction.
  \item[B.] Whatever remains in the raw input buffer from a previous read.
  \item[C.] The current destination before its first block is written.
\end{itemize}

\section{Step 115 Check the complete picture child result}\label{h263-clip-s05}
Invoke the picture encoder and test its status before any file write. A failed
child maps to clip status four. A successful byte count must be nonzero and within
the derived 105\,070-byte picture bound. Retain the returned count across nested
write-helper calls; neither an overwritten return register nor caller capacity
is a substitute for actual encoded length.

\paragraph{Observe} A child returning success with byte count zero does not describe a complete picture.
\Needspace{8\baselineskip}
\paragraph{Check} Which count is written to the stream?
\begin{itemize}
  \item[A.] The entire 131072-byte capacity.
  \item[B.] The admitted complete child byte count.
  \item[C.] The size of all three raw planes.
\end{itemize}

\section{Step 116 Write the encoded bytes with guarded short-I/O handling}\label{h263-clip-s06}
Stage a bounded chunk and its guard byte, then invoke the write syscall with only
the real chunk length. A positive short result advances the original source by
that count, leaving the suffix for another attempt. Reject nonpositive progress.
Keep the accepted-picture record unchanged until both required outputs finish.

\paragraph{Observe} The scratch guard is readable by the syscall but never part of the requested file payload.
\Needspace{8\baselineskip}
\paragraph{Check} What must remain pending after a short write?
\begin{itemize}
  \item[A.] A duplicate of the already accepted prefix.
  \item[B.] Only the next picture header.
  \item[C.] The unaccepted suffix of the current output.
\end{itemize}

\section{Step 117 Accept reconstruction before swapping buffers}\label{h263-clip-s07}
After the stream write succeeds, write all 38\,016 reconstruction bytes with the
same exact-transfer rule. Only after that succeeds promote this destination to
accepted reference, select the other buffer for the next destination, increment
accepted pictures and add the encoded byte count to the accepted total. Publish
those values in the progress record.

\paragraph{Observe} A reconstruction-write failure can leave a second encoded picture in the file while only one picture is accepted.
\Needspace{8\baselineskip}
\paragraph{Check} When does the destination become the reference?
\begin{itemize}
  \item[A.] After both complete output writes are accepted.
  \item[B.] After its first row is reconstructed.
  \item[C.] Immediately after a stream open succeeds.
\end{itemize}

\section{Step 118 Require exact EOF and close owned files}\label{h263-clip-s08}
Repeat until the requested picture count is accepted, then probe one extra input
byte. Accept only zero as exact EOF. A positive or negative result reports input
failure while retaining accepted progress. Close all owned descriptors on this
path and on every error path, exactly once each. Do not infer final success merely
from a progress count that happens to equal the requested count.

\paragraph{Observe} An EOF-probe error after two accepted pictures returns input failure with progress count two.
\Needspace{8\baselineskip}
\paragraph{Check} What establishes overall success?
\begin{itemize}
  \item[A.] A nonempty stream file alone.
  \item[B.] All required work, exact EOF and successful return status zero.
  \item[C.] Any progress value greater than zero.
\end{itemize}

\section{Step 119 Exercise real file errors without host files}\label{h263-clip-s09}
Run the 74 public cases. The grader supplies isolated in-memory files with short
reads/writes, missing paths, descriptor exhaustion and independently targeted
output failures. It decodes every accepted picture and checks levels, motion and
reconstruction using the actual accepted reference. Failure cases check progress,
reference retention and descriptor cleanup. The path-specific failure option
requires the MARS grader revision identified in this package.

\paragraph{Observe} A later encoded-write failure must not be hidden by a previous reconstruction write failing first.
\Needspace{8\baselineskip}
\paragraph{Check} Why target one output path at a time?
\begin{itemize}
  \item[A.] To make all writes succeed regardless of return counts.
  \item[B.] To require a private descriptor number.
  \item[C.] To exercise each failure boundary independently with the real encoder.
\end{itemize}

\section{Step 120 Run and independently decode the generated clip}\label{h263-clip-s10}
Compose your completed picture dependency, then use the public run helper with
a new output directory. It generates the raw input, retains your source instructions,
executes the real MIPS program and passes the resulting stream/reconstruction to
the independent checker. Inspect the saved report and compare both accurate
implementations. Keep finite guest counts separate from whole-process wall time.

\paragraph{Observe} The guarded public reference driver still produces the same 2212-byte stream; its additional I/O work changes the guest count.
\Needspace{8\baselineskip}
\paragraph{Check} What does the host preparation script generate?
\begin{itemize}
  \item[A.] Raw digital input samples and run controls, not the encoded target solution.
  \item[B.] The completed MIPS encoder implementation.
  \item[C.] An all-INTRA substitute for the submitted program.
\end{itemize}
